\section{Worked terminal FTP session}\label{sec:ftp-terminal}

\subsection{Reference exchange from the lecture}
Figure~\ref{fig:ftp-slide-session} is the supplied lecture's directory-listing example. It separates control replies from listing data. It is a reference illustration, not a screenshot of our program. Our client normally negotiates EPSV before trying PASV, so the exact setup reply can differ while the control/data model remains the same.

\begin{figure}[htbp]
\centering
\includegraphics[width=0.95\linewidth,pagebox=cropbox,page=4]{ftp_lecture_figures.pdf}
\caption{Login, passive listing and QUIT in the supplied slides, PDF page~80~\cite{topic17}. The black listing block is data-connection content; the numbered lines are control replies.}
\label{fig:ftp-slide-session}
\end{figure}

\subsection{Run our server and client}
Run both terminals from \code{Group\_Project/FTP}. The example password below is a lab credential, not a saved personal password.

\begin{lstlisting}[style=ftpterm,caption={Terminal 1: start the FTP server.}]
FTP_PASSWORD=lab-demo python3 -m ftp_lab.server \
  --host 127.0.0.1 --port 2121 --root server-data
\end{lstlisting}

\begin{lstlisting}[style=ftpterm,caption={Terminal 2: prepare a file and start the traced client.}]
mkdir -p downloads
printf 'Hello from terminal FTP.\n' > downloads/hello.txt
python3 -m ftp_lab.client --host 127.0.0.1 --port 2121 --trace
\end{lstlisting}
Enter the same password at the client prompt. A trusted-LAN demonstration uses the host's LAN IPv4 address as the server's \code{--host}, and that address as the client's \code{--host} on the other computer. An all-interface server bind is also available with \code{--host 0.0.0.0}. Passive mode needs an incoming negotiated data connection at the FTP server in addition to the control port. No router or firewall rules are changed by these programs.

\begin{lstlisting}[style=ftpterm,caption={Commands entered at the client's ftp prompt.}]
pwd
mkdir experiment
cd experiment
put downloads/hello.txt "hello world.txt"
ls
size "hello world.txt"
get "hello world.txt" downloads/passive-copy.txt
active
get "hello world.txt" downloads/active-copy.txt
passive
ascii
get "hello world.txt" downloads/ascii-copy.txt
binary
delete "hello world.txt"
cdup
rmdir experiment
quit
\end{lstlisting}

Back in the local shell, compare every downloaded copy with the source:
\begin{lstlisting}[style=ftpterm]
cmp downloads/hello.txt downloads/passive-copy.txt
cmp downloads/hello.txt downloads/active-copy.txt
cmp downloads/hello.txt downloads/ascii-copy.txt
\end{lstlisting}
For this LF-terminated text sample, successful commands produce byte-identical copies. Directory creation/removal applies to the server's virtual root; the \code{downloads/} paths refer to the client's own computer.

\subsection{Fresh executed evidence}
On 16 November 2026 we ran an isolated server on a dynamically assigned loopback port and drove the actual terminal client in batch mode. The binary source was all 256 byte values repeated 16 times (4,096 bytes). A separate 34-byte text source included LF and a literal CR. The script tested upload, listing, SIZE, passive and active binary download, ASCII upload/download, remote cleanup, and QUIT.

\lstinputlisting[style=ftpterm,firstline=15,lastline=43,caption={Actual captured control trace: upload, listing, SIZE and passive download. Ephemeral port values are from this run.}]{evidence/ftp_terminal_demo.txt}

\begin{table}[htbp]
\centering\small
\caption{Observed outcomes from the fresh terminal run, not inferred from a success message.}
\begin{tabularx}{\linewidth}{@{}l Y@{}}
\toprule
Check & Observation \\
\midrule
Passive binary copy & Exactly matches the 4,096-byte source. \\
Active binary copy & Exactly matches the same source. \\
ASCII round trip & Exactly matches the 34-byte local text representation. \\
Remote cleanup & Test directory and its files removed; isolated FTP root empty. \\
Trace handling & PASS value appears only as \code{<redacted>}. \\
Client exit & Status 0 after the full batch script. \\
\bottomrule
\end{tabularx}
\end{table}

The accompanying \code{evidence/} directory contains the complete records:
\begin{itemize}[nosep]
\item \path{evidence/ftp_terminal_commands.txt}: batch command input.
\item \path{evidence/ftp_terminal_demo.txt}: chronological terminal trace.
\item \path{evidence/ftp_terminal_verification.json}: sizes, hashes and outcomes.
\end{itemize}
This verification is a real local TCP run; it does not claim that a physical second computer or a particular LAN/firewall path was tested.
